\section{Comparison with base studies and the state of the art}\label{sec:comparison}

The base studies of this work are the 3W dataset papers \citep{vargas2019,vargas2026}, which define the data
and the event classes; the deep-learning detection system of \citet{aranha2024}, whose LSTM autoencoder is
the reference anomaly detector on oil-well sensors; the hybrid Transformer--LSTM of \citet{lai2026}, the
closest forecasting study of a well event with a lead time; PatchTST \citep{nie2023}, on which the sequence
branch of PI-PatchSurv is built; and the open-world detection line of \citet{lopes2025,lopes2026}, which
reports results for choke scaling on 3W. Table~\ref{tab:basepapers} summarises these studies.

\begin{table}[t]
\caption{Base studies: problem, method, data, reported results and main limitations}
\label{tab:basepapers}
\scriptsize\setlength{\tabcolsep}{3pt}
\begin{tabular*}{\textwidth}{@{\extracolsep\fill}p{1.8cm}p{1.7cm}p{2.0cm}p{1.6cm}p{2.2cm}p{2.1cm}}
\toprule
Study & Problem & Method & Data & Reported result & Main limitation \\
\midrule
\citet{vargas2019,vargas2026} & Undesirable-event detection & Dataset and benchmarks & 3W v1 and v2 & Benchmark definitions & No forecasting task \\
\citet{aranha2024} & Anomaly detection & LSTM autoencoder with decision diagram & Field wells, simulated and real & True-positive rate $>$90\,\% in early stages$^*$ & Detection only; split not stated \\
\citet{lopes2025,lopes2026} & Open-world detection and explanation & Autoencoder, classifiers, clustering, LLM agent & Completions; 3W & Novelty detection of choke scaling 94.9\,\% [90.5, 97.3]$^*$ & Detection after onset; not split by well \\
\citet{lai2026} & Timing of water breakthrough & Transformer--LSTM & 9 gas wells, daily data & Breakthrough 2--6 days ahead$^*$ & Other event and field \\
\citet{nie2023} & Long-term forecasting & PatchTST & Public benchmarks & State-of-the-art errors$^*$ & No censoring, no physics \\
\citet{cheng2022} & Scale detection and magnitude & CNN-LSTM-MLP & Simulated histories & $>$95\,\% detection accuracy$^*$ & Simulated data, no timing \\
\botrule
\end{tabular*}
\footnotetext{$^*$Reported in the original study (different task, data or split); not reproduced here.}
\end{table}

\subsection{Quantitative comparison}

The published numbers are not directly comparable with those of this study, because the base studies
address detection or a different event and use different data or splits. To obtain a comparison on the same
data, metrics and protocol, the base architectures were re-implemented (B7 Transformer--LSTM, B8 PatchTST and
B9 LSTM-autoencoder detector) and evaluated exactly as PI-PatchSurv (Table~\ref{tab:quant}).

\begin{table}[t]
\caption{Same-protocol comparison with re-implemented base architectures (leave-one-well-out, 24 events;
mean $\pm$ standard deviation over seeds where several exist). The best deployed baseline is added for
reference}
\label{tab:quant}
\footnotesize\setlength{\tabcolsep}{2.5pt}
\begin{tabular*}{\textwidth}{@{\extracolsep\fill}lccccc}
\toprule
Model & AUPRC@24h & AUROC@24h & Hits/24 & PICP & $\Delta T$ (h) \\
\midrule
B9 LSTM-AE, \citeauthor{aranha2024}$^\dagger$ & 0.262 & 0.365 & 0 & 0.81 & 58.7 \\
B8 PatchTST, \citeauthor{nie2023}$^\dagger$ & $0.569\pm0.041$ & $0.807\pm0.032$ & $2.0\pm1.0$ & 0.52 & 8.8 \\
B7 Transf.--LSTM, \citeauthor{lai2026}$^\dagger$ & $0.711\pm0.031$ & $0.903\pm0.058$ & $13.7\pm0.6$ & 0.33 & 19.6 \\
B4 GRU (best deployed baseline) & $0.712\pm0.123$ & $0.905\pm0.061$ & $14.7\pm1.5$ & 0.43 & 24.2 \\
PI-PatchSurv (proposed) & $0.644\pm0.055$ & $0.794\pm0.018$ & $5.8\pm2.9$ & \textbf{0.87} & \textbf{6.8} \\
\midrule
\citet{lopes2026}, novelty detection$^*$ & \multicolumn{5}{l}{94.9\,\% detection of choke scaling (after onset, not split by well)} \\
\citet{aranha2024}, anomaly detection$^*$ & \multicolumn{5}{l}{True-positive rate $>$90\,\% in early stages (field data)} \\
\botrule
\end{tabular*}
\footnotetext{PICP: coverage of the 90\,\% interval. $^\dagger$Re-implemented and evaluated under the protocol of this study. $^*$Taken from the original study; different task and split, shown for context only.}
\end{table}

Against the re-implemented autoencoder detector, PI-PatchSurv improves AUPRC@24h by 0.382 (146\,\% relative)
and AUROC by 0.429; the gain comes from changing the objective from reconstruction to time-to-event.
Against PatchTST it improves AUPRC by 0.075 (13\,\%), catches 3.8 more events and raises interval coverage
from 0.52 to 0.87; the ablation study attributes this gain to the physics branch and the survival head.
Against the Transformer--LSTM it is lower in AUPRC (by 0.067, or 9\,\%), in AUROC (by 0.109) and in events
caught (by 7.9), but it reduces the time-to-event error from 19.6 to 6.8~h (65\,\%) and raises coverage from
0.33 to 0.87. None of the differences among the deep forecasters is statistically significant
(Section~\ref{sec:significance}).

\subsection{Methodological comparison}

The base detectors model normal operation and flag deviations from it. They answer whether an event is
present, not when it becomes critical, and B9 confirms that a reconstruction score carries almost no
information about the time to the critical level (AUROC 0.37). PatchTST and the Transformer--LSTM forecast
continuous series, but neither handles censoring nor uses process physics. PI-PatchSurv differs from them
in its objective, a censoring-aware survival likelihood with a monotone risk curve; in its inputs, a
self-referenced restriction proxy and physics trends fused through a gate; in its training, which adds
synthetic instances, a monotonicity loss and a deep ensemble; and in its deployment, which uses
self-referenced thresholds and conformal intervals calibrated on separate wells.

\subsection{Fairness of the comparison}

Five factors limit the comparison. First, the published numbers come from detection tasks, from other events
such as water breakthrough or from other data such as completion faults and simulated scale histories, and
most use random or sample-level splits; a held-out-well protocol on 3W lowers macro-F1 from 0.99 to 0.30
\citep{issue231}. Second, the base architectures were re-implemented with the inputs, budget and early
stopping of this study rather than their original hyper-parameters, which Optuna tuning under an equal
budget partly offsets (Appendix~\ref{app:extra}). Third, several publisher pages were not accessible, so only
results stated in accessible abstracts or preprints are quoted. Fourth, seed 0 and seeds 1--2 were trained
on different hardware (Apple M4 and NVIDIA T4), which adds a small amount of variance. Fifth, PI-PatchSurv
does not outperform the best recurrent baselines in ranking or in deployed recall; its advantages are the
quality of its uncertainty, the accuracy of its time-to-event estimate and its robustness to missing sensors.
